\section*{Part V: Measuring Customer Success Rather Than Company Activity}

\subsection*{41. Customer Success Starts Before the Customer Pays}

An organization should define success at the moment it formulates the offer. If a training service sells ``AI awareness,'' the buyer might reasonably expect informed understanding and safer decision-making. If it sells ``production AI capability,'' a higher standard is implied. If it sells measurable workflow improvement, the buyer expects a change in work, not merely a course certificate.

Customer success begins with appropriate selection, explicit prerequisites, a truthful promise, and a baseline. A service can fail the customer by accepting someone it cannot reasonably help, even if every class afterward is competently delivered. The founder must treat \emph{fit} and \emph{outcome conditions} as part of the product, not post-purchase excuses.

\subsection*{42. The Customer and Payer May Disagree About Success}

In education the learner, sponsor, employer, faculty member, and governing institution may want different things. A learner wants usable skill and agency. An employer wants more reliable work and potentially a financial return. A college wants academic quality, compliance, scheduling fit, and student opportunity. A parent or funder may care about cost and future mobility. These goals overlap but are not identical.

Create a stakeholder map with each party's intended benefit, controllable actions, evidence, and possible conflicts. A college may want high pass rates while a serious external assessor refuses to certify unsupported competence. An enterprise client may want rapid automation while workers reasonably fear job redesign. A credible provider should make these tensions explicit rather than designing metrics that flatter the most powerful party.

\subsection*{43. The Five Moments of a Customer Outcome}

For many learning-intensive services, a useful timeline spans 5 distinct moments:

\begin{itemize}[leftmargin=1.5em, itemsep=0.4em]
    \item \textbf{Before:} baseline ability, current workaround, motivation, constraints, and competing options.
    \item \textbf{During:} whether the service was provided, used, understood, and revised in response to feedback.
    \item \textbf{At completion:} independent or authentic demonstration of what was promised.
    \item \textbf{After a delay:} retention and continued use of the acquired capability.
    \item \textbf{In the destination setting:} whether the skill changed actual work, choices, or opportunities under real constraints.
\end{itemize}

Different claims belong to different moments. A course may know a great deal about the 1st 3 and little about the 5th. That is acceptable if disclosed. The problem is advertising the 5th while measuring only the 2nd.

\subsection*{44. Activation, Engagement, Learning, and Success Are Different Variables}

A software user might create an account, log in weekly, and still fail to complete the intended task. A learner might attend every class and remain dependent on tutorials. A consulting client might enjoy weekly updates without ever implementing recommendations.

Measure activity because it can help diagnose where the mechanism breaks. But do not equate it with value. In a subscription business, retention is often central; in a 1-time technical course, appropriate measures include sustained participation, assessed competence, learner regret, consented referrals, and the ability to recruit subsequent cohorts at viable economics. The right indicator depends on the product's economic and functional logic.

\subsection*{45. A Simple Customer-Success Scorecard}

For a technical education service, a small scorecard can distinguish 8 domains:

\begin{table*}[t]
\centering
\small
\caption{The Multi-Domain Customer Success Scorecard}
\begin{tabularx}{\textwidth}{p{0.16\textwidth}p{0.38\textwidth}X}
\toprule
\textbf{Domain} & \textbf{Example Indicator} & \textbf{Main Interpretation Risk} \\
\midrule
Fit & Eligible entrants and baseline distribution & Selectivity can artificially improve final results \\
Experience & Response times, perceived fairness, complaints & Satisfaction is not competence \\
Engagement & Meaningful task attempts, revisions, attendance & Activity is not learning \\
Capability & Independent unseen tasks and rubric results & Narrow sampling and grading variance \\
Durability & Delayed application & Follow-up nonresponse \\
Real-world use & Consented authentic work examples & Context and outside support confound attribution \\
Economics & Contribution, reviewer hours, refund and acquisition cost & Founder time may be underpriced \\
Integrity & Appeals, corrections, overstated claims, privacy incidents & Low complaint counts may reflect fear or silence \\
\bottomrule
\end{tabularx}
\end{table*}

The purpose is not to publish every row publicly. It is to prevent the management team from mistaking strength in 1 area for comprehensive success.

\subsection*{46. The NPS and Satisfaction Trap}

``Would you recommend us?'' can be informative about a customer's attitude, but recommendation likelihood is not verified learning, future referral behavior, or a complete measure of loyalty. A delighted customer may never recommend because they have no relevant peers. A demanding course may produce lower immediate pleasure and stronger retained competence. Conversely, harshness and poor communication are not evidence of rigor.

Treat Net Promoter Score or similar measures as \textbf{experience signals}, inspect free-text reasons, and connect claims to actual referrals or repeat purchase when that is relevant. Do not manipulate students into high ratings by tying grades, certificates, or mentor access to reviews. Solicitation and publication should be voluntary and transparent.
